% Source chapter: nonvf (chapters/nonvf.tex, "Scoring agents for what is neither verifiable nor falsifiable"); moved verbatim from its Introduction.
% Needs the chapter-local \Prob override from macros/nonvf.tex (\Prob{arg} -> \mathbf{P}[arg]); refs nonvf:sec:framework, nonvf:sec:conclusion, nonvf:app:garrabrant resolve to chapter nonvf.
\subsection{Related work}
\label{nonvf:sec:related}

\textbf{Logical uncertainty.} Our work may equivalently be understood as a formalism for \emph{logical uncertainty} (see \citep{halpernReasoningUncertaintySecond2017} for a summary of work in the area), i.e. the problem of assigning probabilities to mathematical sentences in some language $\mathcal{L}$ (which could be a language of FOL sentences). The most relevant work in this domain is \emph{Garrabrant induction} \citep{garrabrantLogicalInduction2020}, which constructs a prediction market for each mathematical sentence that pays off when it is proven by a theorem-enumerator for some theory $\mathcal{T}$ over $\mathcal{L}$.\footnote{to avoid simply being a prediction market for the provability of a sentence, they use a special accounting measure for agent wealths based on ``propositionally consistent worlds'' -- e.g. fixing the combined price of $P$ and $\lnot P$ at \$1 -- which allows them to extend this framework to sentences independent of $\mathcal{T}$.} \draft{A detailed description of the Garrabrant induction construction, and of its limitations for our purposes, is given in Appendix~\ref{nonvf:app:garrabrant}.}

There are two limitations of this framework for our purposes. First: while mathematical theories can easily prove non-VF sentences (e.g. $\forall x,\exists y>x,\,\texttt{prime}(y)$) there is no analogous theory in the real world that we are aware of (``observation'' only proves VF sentences, as discussed). Second: it is philosophically unsatisfying to have probabilities for logical sentences be completely dependent on a mathematical theory, and change when the theory is strengthened or weakened (e.g. by the addition of a G\"odel sentence).

Logical uncertainty has also been discussed in the Bounded Rationality literature e.g. in \citet{halpernAlgorithmicRationalityGame2014,halpernDontWantThink2011,oesterheldTheoryBoundedInductive2023}; however this has typically focused on toy examples based on VF sentences, e.g. betting on the trillionth digit of $\pi$, etc., % (L) closing clause from the long version's corresponding paragraph:
as the games therein have to pay out a reward at a fixed point depending on the answer to the question.

\textbf{Game semantics.} The \emph{verification-falsification game} we use is from Hintikka's \emph{game semantics} \citep{hintikkaGameTheoreticSemantics1997}, which we will discuss in Section~\ref{nonvf:sec:framework}. Game semantics has since been expanded into the field of \emph{computability logic} \citep{japaridzeBeginningWasGame2009,japaridzeSurveyComputabilityLogic2015}, which may be the most natural setting to formulate some of our more speculative ideas about probability theory, described in Sec~\ref{nonvf:sec:conclusion}.

\textbf{Debate.} AI Debate \citep{irvingAISafetyDebate2018} considers sentences $P$ that can be expressed like $\exists x_1\ \forall x_2\ \dots \exists x_n,H(x_1,\dots x_n)$ (more generally any $\Sigma_n$ or $\Pi_n$ form in $H$) where $H$ is some ``human-computable function'', and argues that AIs can be trained to answer such questions correctly via \emph{debate} judged by a human\footnote{in particular, supervised learning can train an AI to learn to answer $\Delta_0$ questions, and reinforcement learning can train an AI to learn to answer $\Sigma_1$ questions}. This is quite similar to our setting, and debate is analogous to the verification-falsification game, except that physical ground truth is replaced by human-computability (or polynomial-computability as they hypothesize). We hope our work can serve as a valuable intuition pump for debate.

\textbf{AI Forecasting.} There has recently been a surge of interest in AI forecasters \citep{zouForecastingFutureWorld2022, schoeneggerLargeLanguageModel2023,yanAutoCastEnhancingWorld2023, halawiApproachingHumanLevelForecasting2024a, schoeneggerWisdomSiliconCrowd2024} -- in particular, recent works such as \citet{sudhirConsistencyChecksLanguage2024} have investigated the \emph{propositional consistency} of AI forecasters, i.e. whether the probabilities they give satisfy basic rules of probability theory like $\Prob{P}+\Prob{\lnot P}=1$. Similarly our work can be understood as asking how we can incentivize or evaluate forecasters respecting the ``continuity of probability'' rules $\Prob{\bigcup P_i}=\lim_{n\to\infty}\bigcup_{i<n}\Prob{P_i}$ and $\Prob{\bigcap P_i}=\lim_{n\to\infty}\bigcap_{i<n}\Prob{P_i}$.
